% ============================================================================
% figures/rag_pipeline.tex
% Pipeline RAG en 6 étapes. Layout linéaire, coordonnées manuelles.
% ============================================================================
\begin{figure}[H]
\centering
\begin{tikzpicture}[
  font=\footnotesize\sffamily,
  every node/.style={align=center, rounded corners=3pt, inner sep=5pt},
  pipebox/.style={draw=blue!40!black, fill=blue!8, thick, minimum width=2.7cm, minimum height=1.3cm},
  arr/.style={-Stealth, thick, gray!70!black},
]

% Étapes (grille horizontale 3 x 2)
\node[pipebox] (s1) at (0, 2) {\textbf{1.\ Ingestion}\\\scriptsize data\_loader.py\\\scriptsize 6 sources};
\node[pipebox] (s2) at (3.4, 2) {\textbf{2.\ Chunking}\\\scriptsize Recursive splitter\\\scriptsize size 500 / overlap 50};
\node[pipebox] (s3) at (6.8, 2) {\textbf{3.\ Embedding}\\\scriptsize MiniLM-L12-v2\\\scriptsize 384 dim., local};

\node[pipebox] (s4) at (0, -0.3) {\textbf{4.\ Indexation}\\\scriptsize ChromaDB persistant\\\scriptsize 9.4 MB};
\node[pipebox] (s5) at (3.4, -0.3) {\textbf{5.\ Retrieval}\\\scriptsize similarité cosinus\\\scriptsize top-$k$ ($k{=}5$)};
\node[pipebox] (s6) at (6.8, -0.3) {\textbf{6.\ Generation}\\\scriptsize GPT-3.5-turbo\\\scriptsize SSE streaming};

% Flèches
\draw[arr] (s1) -- (s2);
\draw[arr] (s2) -- (s3);
\draw[arr] (s3.south) to[bend left=20] (s4.north);
\draw[arr] (s4) -- (s5);
\draw[arr] (s5) -- (s6);

% Phase labels
\node[font=\scriptsize\itshape, text=blue!40!black] at (-1.9, 2) {\rotatebox{90}{\textbf{Indexation}}};
\node[font=\scriptsize\itshape, text=blue!40!black] at (-1.9, -0.3) {\rotatebox{90}{\textbf{Inférence}}};

\end{tikzpicture}
\caption{Pipeline RAG en six étapes. Les trois premières (ingestion, \textit{chunking}, \textit{embedding}) sont exécutées une fois lors de la construction de l'index ; les trois suivantes (indexation, \textit{retrieval}, génération) s'exécutent à chaque requête utilisateur.}
\label{fig:rag_pipeline}
\end{figure}
